\documentclass[11pt]{article}
\usepackage[margin=1in]{geometry}
\usepackage{times}
\usepackage{hyperref}
\hypersetup{colorlinks=true, linkcolor=black, urlcolor=blue, citecolor=black}
\title{The Center of Mass Does Not Negotiate: The Dean Drive and the Recurring Appeal of Oscillating-Mass Propulsion}
\author{Flyxion \\ \small Independent Researcher}
\date{}
\begin{document}
\maketitle

\begin{abstract}
The Dean Drive, patented by Norman Dean in the late 1950s and periodically reinvented under other names since, proposes that an arrangement of eccentrically rotating or vibrating masses within a closed system can produce net unidirectional thrust without expelling any propellant. This paper argues that the device belongs to a broad and recurring family of internal-mass propulsion claims whose common flaw is not a subtle error in a specific gearing scheme but a direct violation of the conservation of linear momentum for an isolated system, a symmetry with no exception carved out for asymmetric timing, mechanical rectification, or any of the variations these devices have proposed across seven decades of reinvention. The paper traces the history of the device through its science fiction publicity, explains the one genuine physical effect that makes such devices appear to work on a table, namely asymmetric friction, and identifies the test that separates the two readings. It makes no claim to have built or tested any such device; its argument is about what the published record permits.
\end{abstract}

\section{Introduction}

Norman Dean's device used a system of rotating eccentric weights coupled through a mechanism intended to apply a stronger force during one phase of the rotation than the opposite phase, producing, in Dean's account, a net thrust in a fixed direction relative to the device housing, with no reaction mass expelled. Dean was granted a patent in 1959 for a system for converting rotary motion into unidirectional motion and another in 1965 for a variable oscillator system, and these two were the only occasions on which he published design details with an explanation, details that are not complete enough to build a working copy \cite{D1959,D1965}. The device became famous outside engineering when the editor John W. Campbell took it up in the pages of \emph{Astounding Science Fiction} around 1960 and argued that it showed the limits of established mechanics \cite{C1960}. Dean claimed demonstrations of the device lifting its own weight or reducing effective weight during operation, though these demonstrations were never conducted under independently controlled conditions. Numerous descendants of the same basic idea, internal masses moved in some cyclic, asymmetric pattern intended to produce net motion of the whole system, have appeared under different names in the decades since, including various self-described inertial propulsion and reactionless drive patents.

\section{Why Internal Asymmetry Cannot Do What Is Claimed}

For a closed system with no external forces acting on it, the center of mass of the system cannot accelerate, regardless of how the masses within the system move relative to each other. This is a direct consequence of Newton's third law applied to the internal forces between the system's components: every internal force has an equal and opposite reaction internal to the same system, so the vector sum of all internal forces is identically zero, and the system's total momentum, and therefore the motion of its center of mass, cannot change. This holds regardless of how asymmetric the timing of the internal motion is, how eccentric the rotating masses are, or how cleverly the mechanical linkage is arranged to make one phase of the cycle look different from another to a casual observer. A device that appears to move forward during one part of its cycle and only partially back during the other is moving its housing relative to an internal counterweight in a way that leaves the system's overall center of mass exactly where conservation of momentum requires it to remain, with any apparent net displacement attributable to the device's mounting, friction with the supporting surface, or measurement over an incomplete cycle.

The theorem has the same status here as in the other reactionless thrust cases in this series. It follows from the translational symmetry of space \cite{N1918}, and it is not the kind of generalization that a single apparatus can overturn. The practical consequence is that the argument against the Dean Drive does not depend on inspecting its mechanism. It would apply equally to a device whose mechanism was never disclosed, which describes much of what Dean built after the patents.

\section{Why Friction Is the Quiet Explanation Every Time}

The reason variations on this idea recur rather than being permanently settled by the momentum argument is that a device resting on a surface with static and kinetic friction is not, technically, a closed system isolated from external forces, since the supporting surface can exert horizontal frictional forces during part of the device's cycle and not during another, if the internal masses' acceleration profile changes the normal force against the surface asymmetrically across the cycle. A vibrating device with the right internal motion can therefore creep across a table, in the manner of a vibrating toy or a stick-slip mechanism, with no violation of anything. The creeping is real motion and it is easy to photograph, and a film of it is indistinguishable, to a viewer who has not tested the device in isolation, from a film of a reactionless drive.

The distinguishing test is therefore simple, and it has been known since the device was first discussed: remove the friction. A suspended device, a device on a low-friction air track, a device on a ballistic pendulum, or a device in free fall cannot use the surface, and the claim of net thrust has to survive these conditions to count. Standard accounts of the episode state that a ballistic pendulum test would have settled the matter and that Dean declined to submit his original device to one \cite{WD}. The point here is not Dean's character, which is not at issue, but the structure of the evidence: a claim whose own proponent withholds the one test that would discriminate the explanations has not met the standard for the claim.

\section{The Family Resemblance}

Seen from outside, the history of such devices is repetitive in a way that is itself informative. Millis, who ran NASA's Breakthrough Propulsion Physics Project, set out the challenge of a space drive that dispenses with reaction mass and placed momentum conservation at the center of the difficulty \cite{M1997}, and the volume that later collected the field's research treats the same problem in the same terms \cite{MD2009}. The inventors differ, the gearings differ, and the demonstrations differ. The claim is the same.

The persistence has a recognizable cause. A rotating eccentric mass produces a force on its housing that points in a different direction at each instant, and an engineer who looks at the force at the instant of maximum asymmetry sees a strong unidirectional pulse. The average over the cycle vanishes, and it is the average that determines motion of the center of mass. A mechanism that is asymmetric within the cycle is a mechanism that fools the eye at one phase and fails the sum.

\section{The Entropic Gravity Framing}

This family of claims does not typically invoke gravitation directly, proposing instead a violation of momentum conservation through mechanical means, but the versions that do claim an accompanying weight reduction effect run into the same absence of coupling mechanism addressed elsewhere in this series: an oscillating internal mass, however configured, has no established or theoretically motivated channel into the large-scale ordered-configuration redistribution that produces gravitational curvature, on the entropic account or on the standard general-relativistic one, and the weight-reduction claims attached to some Dean Drive descendants add an unsupported gravitational layer on top of an already-refuted momentum claim rather than offering independent support for it.

\section{Objections}

Proponents sometimes argue that conservation of momentum, while valid in the timescale-averaged sense, might permit a net displacement over an incomplete or specifically truncated operating cycle, since the theorem strictly applies over a full period of the internal motion. This is occasionally true in a narrow technical sense for systems interacting with an external friction or drag force during only part of their cycle, which is precisely the friction-mediated creeping effect described above, and is not the reactionless, friction-independent thrust the original claims describe; no version of this device family has demonstrated net motion once external asymmetric friction and mounting effects are controlled for.

A second objection is that a device which cannot be built from the patent cannot be fairly tested, and that the successful versions are the ones Dean built later and never disclosed. This is a version of the undisclosed recipe defense that recurs throughout the series. It places the failure on the tester, and it has the effect of making the claim impossible to fail.

A third objection is that the editor of a respected science fiction magazine, himself a trained physicist, took the device seriously. The objection is correct as a description of Campbell, who was an enthusiast for several heterodox ideas of the period, but it is an argument from authority and not from measurement \cite{C1960}. The more informative fact is that the device was taken up by a magazine and not by a laboratory.

\section{Conclusion}

The Dean Drive and its many descendants propose a mechanical route around a conservation law with no carved-out exception for internal timing asymmetry, and the demonstrations that have appeared to succeed have, on closer examination, relied on exactly the external friction interaction that the underlying physics predicts would be necessary to produce real motion, which is to say they have worked, when they have worked at all, for reasons entirely outside the mechanism their inventors described. This paper reports no new tests. It argues that the theory settles the question in advance and that the one test that would discriminate between friction and thrust, a frictionless suspension, is the one that the device's own history never supplied.

\begin{thebibliography}{99}

\bibitem{D1959} N. L. Dean, System for Converting Rotary Motion into Unidirectional Motion, U.S. Patent granted May 1959.

\bibitem{D1965} N. L. Dean, Variable Oscillator System, U.S. Patent granted May 1965.

\bibitem{C1960} J. W. Campbell, The Space Drive Problem, \emph{Astounding Science Fiction}, June 1960.

\bibitem{N1918} E. Noether, Invariante Variationsprobleme, \emph{Nachrichten von der Gesellschaft der Wissenschaften zu Goettingen, Mathematisch-Physikalische Klasse} (1918) 235--257.

\bibitem{WD} Dean Drive, Wikipedia, which summarizes the patent history and the account of the declined pendulum test, https://en.wikipedia.org/wiki/Dean\_drive.

\bibitem{M1997} M. G. Millis, Challenge to Create the Space Drive, \emph{Journal of Propulsion and Power} 13 (1997) 577--582.

\bibitem{MD2009} M. G. Millis and E. W. Davis, editors, \emph{Frontiers of Propulsion Science}, American Institute of Aeronautics and Astronautics, Reston, 2009.

\end{thebibliography}

\end{document}
